\section{Enquadramento teórico}
\label{sec:enquadramento}

Esta secção reúne os conceitos necessários para ler a revisão e a proposta: o
funcionamento do MQTT, os mecanismos de segurança que o protocolo e os seus
\emph{brokers} oferecem nativamente, as primitivas criptográficas usadas pelas
soluções analisadas e o modelo de ameaças que serve de referência à comparação.

% ---------------------------------------------------------------------------
\subsection{O protocolo MQTT}
\label{sec:mqtt}

O MQTT é um protocolo cliente--servidor de publicação/subscrição que funciona
sobre um transporte ordenado e fiável, normalmente TCP (porta 1883) ou TLS
(porta 8883)~\cite{oasis-mqtt311}. A Figura~\ref{fig:mqtt-arquitetura} ilustra os
papéis envolvidos. Os clientes não se conhecem entre si: o desacoplamento no
espaço, no tempo e na sincronização é precisamente o que torna o protocolo
escalável, mas também o que concentra a confiança no \emph{broker}.

\begin{figure}[ht]
  \centering
  \begin{tikzpicture}[
      node distance=1.2cm and 2.2cm,
      cli/.style={draw=iscte, fill=fundotab, rounded corners=2pt, minimum width=2.6cm,
                  minimum height=0.8cm, font=\small, align=center},
      brk/.style={draw=iscte, fill=iscte!15, rounded corners=3pt, minimum width=2.8cm,
                  minimum height=2.8cm, font=\small, align=center},
      seta/.style={-{Stealth[length=2.2mm]}, thick, iscte}]
    \node[brk] (b) {\textbf{\emph{Broker}}\\[2pt]{\footnotesize filtros, ACL,}\\
                     {\footnotesize mensagens retidas}};
    \node[cli, left=of b, yshift=0.9cm] (s1) {Sensor 1\\\footnotesize\emph{publisher}};
    \node[cli, left=of b, yshift=-0.9cm] (s2) {Sensor $n$\\\footnotesize\emph{publisher}};
    \node[cli, right=of b, yshift=0.9cm] (g) {\emph{Gateway}\\\footnotesize\emph{subscriber}};
    \node[cli, right=of b, yshift=-0.9cm] (a) {Atuador\\\footnotesize\emph{subscriber}};
    \draw[seta] (s1) -- node[above, font=\scriptsize, sloped] {PUBLISH \texttt{t}} (b);
    \draw[seta] (s2) -- (b);
    \draw[seta] (b) -- node[above, font=\scriptsize, sloped] {entrega} (g);
    \draw[seta] (b) -- (a);
    \draw[decorate, decoration={brace, amplitude=5pt, mirror}, thick, acento]
      ([yshift=-0.35cm]s2.south west) -- ([yshift=-0.35cm]b.south -| b.west)
      node[midway, below=6pt, font=\scriptsize, text=acento] {TLS (salto 1)};
    \draw[decorate, decoration={brace, amplitude=5pt, mirror}, thick, acento]
      ([yshift=-0.35cm]b.south -| b.east) -- ([yshift=-0.35cm]a.south east)
      node[midway, below=6pt, font=\scriptsize, text=acento] {TLS (salto 2)};
    \draw[decorate, decoration={brace, amplitude=6pt}, thick, bom]
      ([yshift=0.35cm]s1.north west) -- ([yshift=0.35cm]g.north east)
      node[midway, above=7pt, font=\scriptsize, text=bom]
      {proteção fim-a-fim: só esta não depende da confiança no \emph{broker}};
  \end{tikzpicture}
  \caption{Arquitetura de publicação/subscrição do MQTT. O TLS protege cada salto
    em separado e termina no \emph{broker}; só a proteção fim-a-fim da mensagem
    sobrevive a um \emph{broker} não confiável.}
  \label{fig:mqtt-arquitetura}
\end{figure}

\paragraph{Pacotes.} O protocolo define catorze tipos de pacote na versão 3.1.1
(CONNECT, CONNACK, PUBLISH, PUBACK, PUBREC, PUBREL, PUBCOMP, SUBSCRIBE, SUBACK,
UNSUBSCRIBE, UNSUBACK, PINGREQ, PINGRESP e DISCONNECT), a que a versão 5.0
acrescenta AUTH~\cite{oasis-mqtt5}. Cada pacote tem um cabeçalho fixo de dois a
cinco octetos, o que explica a eficiência do protocolo em ligações limitadas.

\paragraph{Tópicos e filtros.} Os tópicos são cadeias UTF-8 hierárquicas
separadas por \texttt{/}. As subscrições usam \emph{filtros} com dois carateres
especiais: \texttt{+} corresponde a exatamente um nível e \texttt{\#} a qualquer
número de níveis, apenas no fim do filtro. Tópicos começados por \texttt{\$}
(como \texttt{\$SYS/\dots}, com estatísticas do \emph{broker}) não são abrangidos
por filtros que comecem por um destes carateres. Formalmente, sendo um tópico
$\tau = (\tau_1, \dots, \tau_m)$ e um filtro $f = (f_1, \dots, f_k)$ sequências
de níveis, define-se a correspondência por
\begin{equation}
  \operatorname{corr}(f, \tau) \iff
  \begin{cases}
    \exists j \le \min(k, m+1):\ f_j = \texttt{\#}\ \wedge\
      \forall i < j:\ f_i \in \{\tau_i, \texttt{+}\}, \\
    k = m\ \wedge\ \forall i \le k:\ f_i \in \{\tau_i, \texttt{+}\},
  \end{cases}
  \label{eq:corr}
\end{equation}
com a restrição adicional de que, se $\tau_1$ começa por \texttt{\$}, então
$f_1 \notin \{\texttt{+}, \texttt{\#}\}$. Esta relação é usada pelas listas de
controlo de acessos e, na proposta, pela restrição chave--tópico
(Secção~\ref{sec:proposta}).

\paragraph{Qualidade de serviço.} O MQTT oferece três níveis de entrega: QoS~0
(no máximo uma vez), QoS~1 (pelo menos uma vez, com possíveis duplicados marcados
pela \emph{flag} DUP) e QoS~2 (exatamente uma vez, com um protocolo de quatro
passos). Note-se que os duplicados legítimos do QoS~1 são indistinguíveis, para um
mecanismo anti-repetição, de uma repetição maliciosa --- uma tensão que a
proposta tem de resolver.

\paragraph{Mensagens retidas e sessões.} O \emph{broker} pode reter a última
mensagem de cada tópico e entregá-la a novos subscritores, e pode manter sessões
persistentes entre ligações. Ambos os mecanismos implicam que o \emph{broker}
armazena conteúdo, o que agrava o problema P1.

% ---------------------------------------------------------------------------
\subsection{Segurança nativa do MQTT e dos \emph{brokers}}
\label{sec:seguranca-nativa}

A Tabela~\ref{tab:seguranca-nativa} resume o que o protocolo e os \emph{brokers}
comuns oferecem sem extensões.

\begin{table}[ht]
  \centering
  \caption{Mecanismos de segurança nativos do MQTT e dos \emph{brokers} e respetivas limitações.}
  \label{tab:seguranca-nativa}
  \small
  \begin{tabularx}{\textwidth}{@{}L{3.1cm}L{4.5cm}Y@{}}
    \toprule
    \textbf{Mecanismo} & \textbf{O que oferece} & \textbf{Limitação} \\
    \midrule
    Utilizador e palavra-passe (CONNECT)
      & Autenticação do cliente pelo \emph{broker}
      & Em claro sem TLS; não autentica mensagens individuais nem o \emph{broker} \\
    \addlinespace
    TLS 1.2/1.3~\cite{rfc5246, rfc8446}
      & Confidencialidade e integridade de cada salto; autenticação do servidor
      & Termina no \emph{broker}; custo de estabelecimento em dispositivos
        restritos~\cite{prantl2021performance}; um cliente só controla o seu
        próprio salto~\cite{spielvogel2021tls} \\
    \addlinespace
    TLS mútuo (certificado de cliente)
      & Autenticação forte do cliente; identidade = nome comum do certificado
      & Exige infraestrutura de chave pública e revogação (CRL) \\
    \addlinespace
    Listas de controlo de acessos
      & Autorização por tópico (leitura/escrita) no \emph{broker}
      & Baseiam-se na identidade declarada se não houver autenticação; não
        protegem contra o próprio \emph{broker} \\
    \addlinespace
    Autenticação estendida (AUTH, MQTT 5)
      & Métodos de desafio--resposta definidos pela aplicação
      & Só autenticação da ligação; não protege o conteúdo \\
    \bottomrule
  \end{tabularx}
\end{table}

Nenhum destes mecanismos protege a \emph{mensagem} de ponta a ponta. É neste
espaço que se situa a maior parte das propostas da literatura e a arquitetura
proposta neste trabalho.

% ---------------------------------------------------------------------------
\subsection{Primitivas criptográficas}
\label{sec:primitivas}

\subsubsection{Cifra autenticada com dados associados}

Uma cifra AEAD é um par de algoritmos que garante, numa só operação,
confidencialidade do texto em claro e integridade e autenticidade tanto do texto
como de dados associados que seguem em claro:
\begin{align}
  (C, T) &= \operatorname{AEAD.Enc}(K, N, P, A), \label{eq:aead-enc}\\
  \operatorname{AEAD.Dec}(K, N, C, T, A) &=
    \begin{cases}
      P & \text{se } T \text{ é válido para } (N, C, A), \\
      \bot & \text{caso contrário,}
    \end{cases} \label{eq:aead-dec}
\end{align}
onde $K$ é a chave, $N$ o \emph{nonce}, $P$ o texto em claro, $A$ os dados
associados (AAD), $C$ o criptograma e $T$ a etiqueta de autenticação. O facto de
$A$ ser autenticado sem ser cifrado é o que permite ligar criptograficamente a
uma mensagem metadados que têm de permanecer legíveis, como o cabeçalho ou o
tópico MQTT.

\paragraph{AES-GCM.} No modo GCM~\cite{nist-sp800-38d}, com chave $K$ e
\emph{nonce} de 96 bits, calcula-se a subchave $H = E_K(0^{128})$ e o bloco
inicial $J_0 = N \,\|\, 0^{31} \,\|\, 1$. O texto é cifrado em modo contador a
partir de $J_0 + 1$ e a etiqueta é
\begin{equation}
  T = \operatorname{MSB}_t\!\left(
        \operatorname{GHASH}_H\!\left(A \,\|\, 0^{v} \,\|\, C \,\|\, 0^{u}
        \,\|\, [\operatorname{len}(A)]_{64} \,\|\, [\operatorname{len}(C)]_{64}\right)
        \oplus E_K(J_0)\right),
  \label{eq:gcm}
\end{equation}
onde $\operatorname{GHASH}_H$ é uma função polinomial em $\mathrm{GF}(2^{128})$ e
$u, v$ completam blocos de 128 bits. Com $P$ vazio obtém-se o GMAC, um código de
autenticação de mensagem: $T = \operatorname{AEAD.Enc}(K, N, \varnothing, A)$.

\paragraph{Limite dos \emph{nonces} aleatórios.} A segurança do GCM colapsa se o
mesmo par $(K, N)$ for usado duas vezes: o XOR dos dois criptogramas revela o XOR
dos textos e a subchave $H$ pode ser recuperada, permitindo falsificar
etiquetas~\cite{bock2016nonce}. Com \emph{nonces} aleatórios de $b = 96$ bits e
$q$ mensagens por chave, a probabilidade de colisão é limitada pelo paradoxo do
aniversário:
\begin{equation}
  \Pr[\text{colisão}] \;\le\; \frac{q(q-1)}{2^{b+1}} \;\approx\; \frac{q^{2}}{2^{97}}.
  \label{eq:aniversario}
\end{equation}
Para $q = 2^{32}$ a probabilidade é da ordem de $2^{-33}$, razão pela qual o NIST
limita a $2^{32}$ o número de invocações por chave com \emph{nonces}
aleatórios~\cite{nist-sp800-38d}. Esquemas de longa duração precisam, por isso,
de rotação de chaves.

\paragraph{ChaCha20-Poly1305.} Alternativa sem dependência de instruções AES em
\emph{hardware}, com chave de 256 bits, \emph{nonce} de 96 bits e etiqueta de 128
bits~\cite{rfc8439}; é a escolha de~\cite{sadio2019lightweight} para nós
restritos.

\subsubsection{Códigos de autenticação de mensagem}

O HMAC~\cite{rfc2104} constrói um MAC a partir de uma função de \emph{hash} $\mathcal{H}$:
\begin{equation}
  \operatorname{HMAC}(K, m) =
    \mathcal{H}\big((K' \oplus \mathit{opad}) \,\|\,
    \mathcal{H}((K' \oplus \mathit{ipad}) \,\|\, m)\big),
  \label{eq:hmac}
\end{equation}
onde $K'$ é a chave ajustada ao tamanho do bloco. Um MAC dá integridade e
autenticidade face a quem não conhece $K$, mas \emph{não} prova a origem perante
terceiros: qualquer detentor de $K$ (incluindo o recetor) pode produzir a mesma
etiqueta.

\subsubsection{Assinaturas digitais Ed25519}

O Ed25519~\cite{bernstein2012ed25519, rfc8032} é um esquema de assinatura sobre a
curva de Edwards equivalente à Curve25519, com ponto base $B$ de ordem prima
$\ell$. A partir da chave privada (32 octetos) derivam-se o escalar $s$ e a chave
pública $A = [s]B$. Para assinar $M$:
\begin{equation}
  r = \mathcal{H}(\mathit{prefixo} \,\|\, M) \bmod \ell,\quad
  R = [r]B,\quad
  k = \mathcal{H}(R \,\|\, A \,\|\, M) \bmod \ell,\quad
  S = (r + k\,s) \bmod \ell,
  \label{eq:ed25519-sign}
\end{equation}
e a assinatura é $(R, S)$, com 64 octetos. A verificação aceita sse
\begin{equation}
  [8][S]B = [8]R + [8][k]A.
  \label{eq:ed25519-verify}
\end{equation}
Como $r$ é derivado deterministicamente da chave e da mensagem, o esquema elimina
a dependência de um gerador aleatório no momento da assinatura. Ao contrário de um
MAC, só o detentor de $s$ consegue assinar, enquanto qualquer detentor de $A$
consegue verificar: a assinatura dá \emph{prova de origem} (não repúdio).

% ---------------------------------------------------------------------------
\subsection{Validade temporal e proteção contra repetição}
\label{sec:validade}

\paragraph{Janela temporal.} Uma mensagem tem \emph{validade temporal} se foi
produzida há pouco tempo, segundo um carimbo temporal autenticado pelo remetente.
Dada uma mensagem com carimbo temporal $t_m$,
recebida no instante $t$, um recetor com tolerância de idade $\Delta_a$ e de
desvio de relógio $\Delta_s$ aceita-a sse
\begin{equation}
  t - \Delta_a \;\le\; t_m \;\le\; t + \Delta_s .
  \label{eq:validade}
\end{equation}
A validade temporal responde à pergunta \enquote{esta mensagem é recente?} e
limita o tempo durante o qual uma mensagem capturada é útil, mas \emph{não}
responde a \enquote{esta mensagem já foi recebida?}: não impede a repetição
dentro da janela $\Delta_a$. É esse o papel do mecanismo seguinte.

\paragraph{Janela deslizante de números de sequência.} O mecanismo anti-repetição
do IPsec ESP~\cite{rfc4303} mantém, por remetente, o maior número de sequência
aceite $h$ e um mapa de bits $W \in \{0,1\}^{w}$ (com $w = 64$ o valor
recomendado), em que $W_i = 1$ significa que $h - i$ já foi recebido. Para uma
mensagem autêntica com número $s > 0$:
\begin{equation}
  \operatorname{aceita}(s) \iff
  \begin{cases}
    s > h, & \text{e então } W \leftarrow (W \ll (s - h)) \mid 1,\ \ h \leftarrow s; \\
    0 \le h - s < w \ \wedge\ W_{h-s} = 0, & \text{e então } W_{h-s} \leftarrow 1; \\
    \text{falso}, & \text{se } h - s \ge w \text{ ou } W_{h-s} = 1.
  \end{cases}
  \label{eq:janela}
\end{equation}
O mecanismo tolera reordenação até $w$ posições e rejeita duplicados e mensagens
antigas. Duas condições são essenciais: o estado só pode ser atualizado
\emph{depois} de verificada a autenticidade (caso contrário um atacante avança a
janela com números forjados), e cada chave deve pertencer a um único remetente.

% ---------------------------------------------------------------------------
\subsection{Limitação de tráfego}
\label{sec:token-bucket}

O balde de testemunhos (\emph{token bucket}) limita a taxa média a $\rho$
mensagens por segundo, com rajadas até $\beta$. O número de testemunhos
disponível no instante $t$, dado o estado no último acesso $t_0$, é
\begin{equation}
  b(t) = \min\!\big(\beta,\ b(t_0) + \rho\,(t - t_0)\big),
  \label{eq:token-bucket}
\end{equation}
e uma mensagem é aceite sse $b(t) \ge 1$, caso em que $b \leftarrow b(t) - 1$. Num
intervalo de duração $\Delta$, o número de mensagens aceites é limitado por
$\beta + \rho\,\Delta$, independentemente da taxa oferecida.

% ---------------------------------------------------------------------------
\subsection{Modelo de ameaças}
\label{sec:ameacas}

Considera-se um adversário que pode (i)~observar e alterar o tráfego na rede,
(ii)~ligar-se ao \emph{broker} como cliente, subscrever, publicar e repetir
mensagens, (iii)~reivindicar a identidade de outro cliente, (iv)~comprometer o
próprio \emph{broker} ou (v)~comprometer um dispositivo legítimo, obtendo as suas
chaves. A Tabela~\ref{tab:stride} organiza as ameaças segundo as categorias STRIDE,
que correspondem às classes de problemas identificadas para o MQTT
em~\cite{hintaw2023mqtt}.

\begin{table}[ht]
  \centering
  \caption{Ameaças ao MQTT segundo a classificação STRIDE e mecanismos que as mitigam.}
  \label{tab:stride}
  \small
  \begin{tabularx}{\textwidth}{@{}L{2.4cm}L{5.2cm}Y@{}}
    \toprule
    \textbf{Categoria} & \textbf{Ameaça no MQTT} & \textbf{Mitigação típica} \\
    \midrule
    \emph{Spoofing}    & Cliente reivindica a identidade de outro; mensagens injetadas em tópicos alheios
                       & Autenticação (palavras-passe, mTLS); autenticação da mensagem; restrição chave--tópico \\
    \emph{Tampering}   & Alteração do conteúdo, do cabeçalho ou do tópico em trânsito ou no \emph{broker}
                       & TLS (por salto); AEAD/MAC/assinatura com tópico nos dados associados (fim-a-fim) \\
    \emph{Repudiation} & Remetente nega ter enviado; recetor fabrica mensagens em nome do remetente
                       & Assinaturas digitais (um MAC partilhado não basta) \\
    \emph{Information disclosure} & Subscrição de \texttt{\#}; \emph{broker} curioso; tópicos \texttt{\$SYS}
                       & ACL; cifra fim-a-fim; TEE no \emph{broker}~\cite{segarra2020mqttz} \\
    \emph{Denial of service} & Inundação de mensagens; ligações lentas~\cite{vaccari2020slowite}; esgotamento de estado
                       & Limitação de tráfego; limites de pacotes e ligações; estado limitado; IDS~\cite{vaccari2020mqttset} \\
    \emph{Elevation of privilege} & Abuso de ACL baseadas em identidade não verificada; chave de um dispositivo usada para outros
                       & ACL com identidade autenticada; restrição chave--tópico; revogação \\
    \bottomrule
  \end{tabularx}
\end{table}

\paragraph{Defesa em profundidade.} Se um ataque tiver de ultrapassar $n$ camadas
de defesa independentes, cada uma contornável com probabilidade $p_i$, a
probabilidade de sucesso é
\begin{equation}
  P_{\text{sucesso}} = \prod_{i=1}^{n} p_i ,
  \label{eq:profundidade}
\end{equation}
estritamente menor do que a de qualquer camada isolada sempre que $p_i < 1$. A
hipótese de independência raramente se verifica por inteiro (por exemplo, o roubo
de uma chave de dispositivo pode afetar várias camadas), mas o modelo justifica a
combinação de mecanismos em camadas distintas: transporte, mensagem e
disponibilidade.
